\section{附录三：模型实验室自检评分卡}

为了把本期访谈转成可执行判断，可以给任何模型实验室或 Agent 创业团队做一个简化评分卡。它不是排名工具，而是帮助读者判断一个团队是否真的理解 Agent 范式。

\noindent\begin{longtable}{p{0.24\textwidth}p{0.33\textwidth}p{0.34\textwidth}}
\toprule
维度 & 高分表现 & 低分信号 \\
\midrule
范式判断 & 能清楚说明 Chat、Reasoning、Agent 三阶段的训练和产品差异 & 只把 Agent 当成工具调用或复杂 prompt。 \\
后训练系统 & 有任务环境、rollout、verifier、eval、成本监控的闭环 & 只有离线数据或人工挑样例。 \\
模型编排 & 能按任务阶段路由大小模型、多模态模型和专用模型 & 所有任务都调用同一个最大模型。 \\
算力配置 & research、pre-train、post-train 有明确比例和迭代节奏 & 资源被历史路径锁死，后训练没有足够卡。 \\
组织机制 & 预训练、后训练、产品、系统、评估共同参与决策 & 单一团队闭门决定，真实用户和工程反馈进不来。 \\
成本意识 & 端到端完成率、延迟、价格一起优化 & 只关心能力展示，不关心单位任务经济账。 \\
价值观与边界 & 明确哪些任务应自动化、哪些需要人类保留判断 & 只追求替代和增长，不讨论安全、责任和社会后果。 \\
\bottomrule
\end{longtable}

\begin{knowledgebox}{如何使用这张评分卡}
读后续Zhang Xiaojun AI 队列时，可以把每期嘉宾观点映射到这七个维度。比如苏煜访谈更偏“范式判断”和“长期学习”，罗福莉访谈则更偏“后训练系统”“模型编排”“算力配置”和“组织机制”。这样多期笔记会逐渐形成一个可比较的 AI Lab 操作系统图谱。
\end{knowledgebox}

\subsection{本章小结}

评分卡把访谈里的判断变成可复用框架。它提醒读者：Agent 时代不是某个单点技术胜出，而是多个系统维度同时达标。

